\section{Discussion: Why Behavioral Video Matters}
\label{sec:discussion}

Our results position video as a non-invasive complement to EEG: on this cohort
it matches EEG at detection and surpasses it where behavior \emph{is} the
endpoint. Four reasons make the video route worth taking seriously in
preclinical seizure work.

\noindent\textbf{Detection without electrodes.}
For binary detection the video model reaches 0.968 macro-F1 / 0.993
AUROC, against the EEG GRU's 0.974 / 0.996---a 0.006 margin to EEG that is small
but consistent across five seeds (Welch $p{=}0.003$; Table~\ref{tab:results}).
Detection is thus a \emph{parity} result rather than a win for either modality, and it is architecture-independent
on both sides: it holds across ten video backbones and seven EEG
models (Tables~\ref{tab:results},~\ref{tab:eegarch},~\ref{tab:vidarch}). A
purely behavioral, electrode-free model therefore preserves detection accuracy,
echoing the human vEpiNet system, where video cut false
positives by a third~\cite{vepinet2024}.

\noindent\textbf{The endpoint is behavioral.}
Seizure \emph{severity} is defined behaviorally. The Racine
scale~\cite{racine1972} grades motor manifestations such as twitching, rearing,
rhythmic forelimb clonus, and rearing-and-falling. Video observes the very
signal the clinical grade is built on, so it degrades gracefully as the task
refines (macro-F1 $0.968\!\to\!0.780\!\to\!0.657$). The EEG reference, one step
removed from motor behavior, declines more steeply from a marginally higher start
($0.974\!\to\!0.720\!\to\!0.551$): the two modalities are level at detection and
separate only as the task becomes explicitly behavioral. The gap widens at the
rarest, most severe stage. Convulsive motion is unmistakable on camera yet poorly
separated in a single EEG channel, so at five-class grading the Stage\,5 per-class
F1 reaches 0.42 for video against 0.12 for EEG (Table~\ref{tab:perclass}). This
ordering survives the stricter subject-disjoint protocol, though the margin narrows:
held-out animals cost video more than they cost the EEG reference at five classes
($0.657\!\to\!0.475$ against $0.551\!\to\!0.404$), leaving a gap of 0.071 rather
than 0.106 (Table~\ref{tab:subjectwise}).

\noindent\textbf{Scale, cost, and welfare.}
Electrode EEG requires surgical implantation: expensive, low-throughput, and a
source of infection, signal drift, and animal-welfare burden. A camera is cheap,
surgery-free, and able to watch many cages at once. Video-only screening
therefore scales to the large cohorts preclinical drug discovery needs, and it
reduces harm in the spirit of the 3Rs. Deployment is equally favorable: the
4M-parameter X3D-M matches the 33M R(2+1)D (Tables~\ref{tab:vidarch}
and~\ref{tab:vidspec}), making
real-time, edge-side monitoring realistic.

\noindent\textbf{Interpretability.}
The Grad-CAM maps of Sec.~\ref{sec:application} (Fig.~\ref{fig:gradcam})
confirm that the detector
attends to behavior rather than recording-setup shortcuts, a prerequisite for
trusting a non-invasive screen.

\noindent\textbf{What EEG still provides.}
Electrodes remain indispensable. EEG captures electrographic-only
(non-convulsive) seizures that lack a visible motor correlate, and it recovers
precise onset timing and localization---signals to which a camera is blind. The
two modalities are therefore complementary. The strongest system would fuse
them---the near-empty ``video--EEG, animal'' cell that our study motivates and
that we outline as future work (Sec.~\ref{sec:open}). What our results establish
is that the \emph{behavioral half} of the video--EEG gold standard is far more
automatable than its scarce use suggests.

\begin{takeaway}
\takeawaylabel The behavioral half of the video--EEG gold standard is substantially automatable:
a camera alone matches electrode-based detection, grades severity more gracefully,
and generalizes more stably across animals. Electrodes remain necessary for
electrographic-only events and for precise onset timing.
\end{takeaway}

